\section*{Chapitre \ref{ch:ml:linear-and-regularised-baselines} --- Modèles de référence linéaires et régularisés}

\begin{solution}{exo:ml:linear-and-regularised-baselines:1}
Chaque côté vend 30 positions de $1/50$ et en achète 30~: $0.6 + 0.6 = 1.2$ échangé par côté, 2.4 au total par unité de capital,
soit 1.2 fois l'exposition brute de 2. Coût~: $2.4\times20\,\text{pb} = 0.48~\%$ du capital de ce mois.
\end{solution}

\begin{solution}{exo:ml:linear-and-regularised-baselines:2}
$1.6/2.5\times\sqrt{12} = 2.22$~; après coûts $1.4/2.5\times\sqrt{12} = 1.94$.
\end{solution}

\begin{solution}{exo:ml:linear-and-regularised-baselines:3}
Le long de $(1, 1)/\sqrt2$, la partie commune, qui porte $(1 + 0.9)/2 = 95~\%$ de la variance. C'est la bonne direction
pour une régression seulement si la cible charge sur la partie commune~; si elle charge sur la différence des deux variables, la
dernière composante de l'ACP est celle qui compte, et la \omterm{def:ml:linear-and-regularised-baselines:pcr}{régression sur composantes principales} (PCR) à une composante ne trouve rien.
\end{solution}

\begin{solution}{exo:ml:linear-and-regularised-baselines:4}
Un portefeuille par déciles n'utilise que les extrêmes du classement, là où les deux modèles s'accordent le plus, et il ignore la taille de
des prévisions, là où se loge en partie le meilleur R² des arbres. Les prévisions des arbres bougent aussi davantage (taux de rotation 0.97
contre 0.80), donc elles paient davantage de coûts~: brut 2.20 contre 2.17, net 1.90 contre 1.91.
\end{solution}

\begin{solution}{exo:ml:linear-and-regularised-baselines:5}
Le portefeuille par déciles dépend du classement (l'IC)~; le R² dépend aussi de l'échelle des prévisions. Avec deux
ans de données, les arbres classent assez bien mais leurs prévisions s'étalent trop, plus du double de l'échelle optimale
$c^\star$ de la \cref{prop:ml:why:ceiling}, donc le R² est négatif alors que le classement rapporte un ratio de Sharpe de 1.11.
\end{solution}

\begin{solution}{exo:ml:linear-and-regularised-baselines:6}
Une comparaison n'est valable que sur les mêmes données, découpages, cible, indice de référence (zéro ou la moyenne) et période~: le R² varie
de plusieurs dixièmes de point d'une décennie à l'autre (chapitre~1). Il faut réajuster le \omterm{def:ml:linear-and-regularised-baselines:baseline}{modèle de référence} linéaire dans un banc d'essai propre et
comparer les résultats là-bas, avec l'erreur type de l'écart mensuel.
\end{solution}

\begin{solution}{exo:ml:linear-and-regularised-baselines:7}
\texttt{table(style=0.0)}~: ratios de Sharpe nets de 3.68 pour la régression ridge et 4.86 pour les arbres (4.24 pour le lasso), contre 1.59
et 1.90. Le R² bouge moins (0.35~\% et 0.57~\% contre 0.30~\% et 0.47~\%). Les rendements des facteurs ajoutent une variance que
un portefeuille par déciles, trié sur les caractéristiques qui les portent, ne peut diversifier~; le R² compte la variance de chaque
action, dont le risque facteur n'est qu'une petite part.
\end{solution}

\begin{solution}{exo:ml:linear-and-regularised-baselines:8}
Avec $X^\top X = I$, l'objectif du lasso se sépare~: $\frac12(b_j - \beta_j)^2 + \lambda|\beta_j|$ pour chaque
coefficient des moindres carrés ordinaires $b_j$, minimisé par $\sign(b_j)\max(|b_j| - \lambda, 0)$. Dès que $\lambda$ dépasse le plus grand $|b_j|$, tous les
coefficients sont nuls et le modèle prédit la moyenne (nulle) de la cible centrée, dont le R² par rapport à zéro vaut zéro~: la valeur du chemin en
$\lambda = 3\times10^{-3}$. Les caractéristiques classées sont presque orthogonales, donc la même logique tient approximativement.
\end{solution}

\begin{solution}{pb:ml:linear-and-regularised-baselines:1}
\begin{enumerate}
\item 0.71~\% sur les années de test~; caractéristiques 0 à 5.
\item Pour que les portefeuilles triés sur des caractéristiques portent un risque facteur, comme le font les vrais~; ils donnent des ratios de Sharpe par déciles
allant d'environ 5 à environ 2.
\item 2.45 pour leur meilleur réseau, 0.83 pour un modèle linéaire à trois variables.
\item Les rendements des facteurs font des mois, et non des lignes, l'échantillon effectif (chapitre~1).
\item Ridge $\lambda = 10^4$~; lasso $10^{-3}$ (six coefficients retenus)~; à $3\times10^{-3}$, le lasso n'en retient aucun et obtient
zéro.
\item Les caractéristiques classées sont presque non corrélées et de même variance, donc les composantes principales sont des directions arbitraires
(voir la \cref{prop:ml:baselines:pls})~; le signal est réparti sur beaucoup d'entre elles.
\item Sa première composante est $X^\top y$, la direction des moindres carrés à un facteur d'échelle près lorsque les variables sont presque orthogonales.
\item Le carré et la valeur absolue (non-linéarités additives), et non l'interaction $c_0c_3$.
\item Arbres 0.47~\%, IC 0.076, ratio de Sharpe net 1.90~; lasso 0.32~\%, IC 0.066, ratio de Sharpe net 1.91.
\item Le R² valorise les tailles et toutes les actions~; le portefeuille par déciles valorise les extrêmes du classement, nets de coûts, et les
arbres échangent davantage.
\item Entre cinq et dix ans (60 mois~: $-0.02~\%$ contre 0.18~\%~; 120 mois~: 0.41~\% contre 0.34~\%).
\item R² $-0.71~\%$ mais un ratio de Sharpe net de 1.11 (ridge 0.90)~: le classement est utile, les tailles ne le sont pas.
\item \emph{Résultat nommé}~: ridge 0.30~\% et IC 0.060 contre 0.47~\% et 0.076 pour les arbres (le lasso~: 0.32~\% et 0.066)~;
les arbres dépassent la régression ridge en R² entre cinq et dix ans d'historique, et ne battent jamais le
ratio de Sharpe net du lasso.
\item Le lasso~: performance nette égale, rotation plus faible, six coefficients lisibles et surveillables.
\item Ridge sur de nombreuses variables non linéaires aléatoires construites à partir des caractéristiques, fortement pénalisé.
\item L'écart net mensuel a une moyenne de $-0.08~\%$ et un écart-type de 2.56~\%~: $(2\times2.56/0.08)^2$ est
plus de 4\,000 mois~; aucune trajectoire réalisable ne les sépare.
\item Caractéristiques corrélées par blocs, le signal portant sur la partie commune de chaque bloc.
\item Le ratio de Sharpe net et la rotation, à côté de l'IC et de son $t$.
\item La taille de chaque mouvement, qui est surtout du bruit ici~: la \omterm{def:ml:linear-and-regularised-baselines:logistic}{régression logistique} classe presque aussi bien (IC 0.060).
\item Un \omterm{def:ml:linear-and-regularised-baselines:baseline}{modèle de référence} (baseline) chiffre la complexité~: c'est le nombre qu'un nouveau modèle doit battre, à conditions égales, d'une marge supérieure au bruit.
\end{enumerate}
\end{solution}

\begin{solution}{iq:ml:linear-and-regularised-baselines:1}
Une régression linéaire pénalisée (ridge ou lasso) sur des variables classées par date, la cible étant centrée par date, la pénalité
choisie par validation croisée purgée sur des dates entières, et le rapport standard (R² par rapport à zéro, IC avec $t$ HAC,
portefeuille par déciles net de coûts). Chaque modèle ultérieur se juge par rapport à lui.
\iqlookfor{classement, centrage, découpages purgés par date, et un rapport allant au-delà du R².}
\end{solution}

\begin{solution}{iq:ml:linear-and-regularised-baselines:2}
Il faut essayer les deux~: ridge si de nombreuses caractéristiques portent de petits effets, lasso si quelques-unes en portent l'essentiel~; avec des variables classées, presque orthogonales,
la parcimonie du lasso ne coûte guère et son modèle reste lisible. L'elastic net couvre les groupes corrélés. Il faut choisir
par validation croisée purgée, et non au goût.
\iqlookfor{l'a priori de parcimonie, variables corrélées et choix empirique.}
\end{solution}

\begin{solution}{iq:ml:linear-and-regularised-baselines:3}
La PCR retient les directions de plus forte variance des variables, en ignorant la cible~; les \omterm{def:ml:linear-and-regularised-baselines:pcr}{moindres carrés partiels} (PLS) retiennent celles de plus forte covariance avec
la cible. Elles diffèrent lorsque le signal se situe dans des directions de faible variance, ou lorsque les variables ont des variances similaires (rangs),
où les directions de l'ACP sont arbitraires.
\iqlookfor{réduction supervisée ou non supervisée, et le cas où cela compte.}
\end{solution}

\begin{solution}{iq:ml:linear-and-regularised-baselines:4}
Le R² récompense les bonnes tailles et chaque action~; un portefeuille récompense le classement aux extrêmes et paie la rotation.
Un modèle peut améliorer les tailles au milieu de la distribution, ou échanger davantage, et laisser le portefeuille inchangé, voire pire.
\iqlookfor{l'écart entre précision des prévisions et valeur de trading.}
\end{solution}

\begin{solution}{iq:ml:linear-and-regularised-baselines:5}
Classer les variables par date, centrer la cible par date, construire des découpages purgés de dates entières avec un embargo égal à la longueur de l'étiquette,
évaluer une grille logarithmique assez large pour que ses deux extrémités soient moins bonnes, et retenir la valeur la plus pénalisée à moins d'une
erreur type du meilleur.
\iqlookfor{découpages par date, une grille qui encadre l'optimum, et la règle de l'erreur type.}
\end{solution}

\begin{solution}{iq:ml:linear-and-regularised-baselines:6}
Maximiser $w^\top X^\top y$ pour $w$ unitaire~: par Cauchy--Schwarz, $w\propto X^\top y$. La première composante principale maximise
$w^\top X^\top Xw$ à la place, ce qui ignore $y$~; lorsque la cible dépend d'une direction de faible variance, les \omterm{def:ml:linear-and-regularised-baselines:pcr}{moindres carrés partiels} la trouvent et l'ACP
ne la trouve pas.
\iqlookfor{la dérivation et la raison géométrique.}
\end{solution}
